%------------------------------------------------------------------------------%

\usepackage{calculo_varias_variaveis-1}

\title {Continuidade de Funções de Várias Variáveis}

%------------------------------------------------------------------------------%
\begin{document}

\addtitle

%------------------------------------------------------------------------------%
\section{Continuidade}

%------------------------------------------------------------------------------%
\begin{frame}{Definição}
  Uma função \(f(x, y)\) é \emph{contínua no ponto} $(x_0, y_0)$, se
  \begin{enumerate}[<+->]
    \setlength\itemsep{1em}
    \item $f$ é definida em \((x_0, y_0)\)
    \item \(\lim_{(x, y)\to(x_0, y_0)} f(x, y)\) existe
    \item \(\lim_{(x, y)\to(x_0, y_0)} f(x, y) = f(x_0, y_0)\)
  \end{enumerate}

  \vspace{\baselineskip}
  \onslide<+->{Uma \emph{função é contínua} se for contínua em todos os pontos do seu domínio}
\end{frame}

%------------------------------------------------------------------------------%
\begin{frame}{Limite de Função Contínua}
  Se \(f(x, y)\) é contínua em $(x_0, y_0)$,
  \pause
  então
  \[
    \lim_{(x, y)\to(x_0, y_0)} f(x, y) \pause = f(x_0, y_0)
  \]
\end{frame}

%------------------------------------------------------------------------------%
\begin{frame}{Combinações Algébricas}
  Combinações algébricas de funções contínuas são funções contínuas

  \pause
  \vspace{\baselineskip}
  Aplicação direta das propriedades de limites
\end{frame}

%------------------------------------------------------------------------------%
\addtocounter{exemplo}{1}
\begin{frame}{Exemplo \theexemplo}
  Exemplos de funções contínuas no plano
  \begin{enumerate}[<+->]
    \item \(f(x, y) = x\)
    \item \(f(x, y) = y\)
    \item \(f(x, y) = xy\)
    \item \(f(x, y) = x^n\) para $n$ inteiro positivo
    \item \(f(x, y) = y^n\) para $n$ inteiro positivo
  \end{enumerate}
\end{frame}

%------------------------------------------------------------------------------%
\addtocounter{exemplo}{1}
\begin{frame}{Exemplo \theexemplo}
  Mostre que
  \[
    f(x, y) =
    \begin{cases}
      \dfrac{2xy}{x^2+y^2}, & (x, y) \neq (0, 0) \\
      \quad\; 0,            & (x, y) =    (0, 0)
    \end{cases}
  \]
  é contínua em todos os pontos, exceto na origem.
\end{frame}

%------------------------------------------------------------------------------%
\begin{frame}{Exemplo \theexemplo}
  Consideramos primeiro o caso \((x, y) \neq (0, 0)\)

  \pause
  \vspace{\baselineskip}
  \(2xy\) é uma função contínua no plano

  \pause
  \(x^2+y^2\) é uma função contínua no plano

  \pause
  então \(\frac{2xy}{x^2+y^2}\)
  \pause
  é contínua para todo \((x, y) \neq (0, 0)\)
\end{frame}

%------------------------------------------------------------------------------%
\begin{frame}{Exemplo \theexemplo}
  Em \((x, y) = (0, 0)\) temos que aplicar a definição
  \pause
  \begin{enumerate}
    \item \(f\) é definida em \((0, 0)\)
    \item \(\lim_{(x, y)\to(0, 0)} f(x, y)\) existe
    \item \(\lim_{(x, y)\to(0, 0)} f(x, y) = f(0, 0)\)
  \end{enumerate}

  \pause
  \vspace{\baselineskip}
  A condição 1 vale pois \(f(0, 0) = 0\)
\end{frame}

%------------------------------------------------------------------------------%
\begin{frame}{Exemplo \theexemplo}
  Verificando a condição 2, percebemos que o limite não existe.

  \pause
  Se fizermos \(y=mx\),
  \pause
  para \(m\neq 0\) e \(x\neq 0\), temos
  \pause
  \[
    f(x, y)\evalat{y=mx}{}
    \pause
    = \frac{2xy}{x^2+y^2}\evalat{y=mx}{}
    \pause
    = \frac{2mx^2}{x^2+m^2x^2}
    \pause
    = \frac{2mx^2}{\left(1+m^2\right)x^2}
    \pause
    = \frac{2m}{1+m^2}
  \]

  \pause
  Para cada valor de $m$, a função $f$ se aproxima da origem com um valor diferente.

  \pause
  Portanto, o limite não existe e a função não é contínua na origem
\end{frame}

%------------------------------------------------------------------------------%
\begin{frame}{Teste dos Dois Caminhos}
  Se \(f(x, y)\) tem limites diferentes ao longo de dois caminhos
  diferentes,

  \pause
  no domínio de $f$,

  \pause
  quando $(x, y)$ se aproxima de $(x_0, y_0)$,

  \pause
  então
  \[
    \lim_{(x, y)\to(x_0, y_0)} f(x, y) \quad \text{ \emph{não existe}}
  \]
\end{frame}

%------------------------------------------------------------------------------%
\addtocounter{exemplo}{1}
\begin{frame}{Exemplo \theexemplo}
  Verifique se a função
  \(
    f(x, y) = \frac{2x^2y}{x^4+y^2}
  \)
  possui limite na origem.

  \pause
  \vspace{\baselineskip}
  O limite na origem leva a uma indeterminação do tipo $\sfrac{0}{0}$
\end{frame}

%------------------------------------------------------------------------------%
\begin{frame}{Exemplo \theexemplo}
  Fazendo \(y=mx\), para \(m\neq 0\) e \(x\neq 0\),
  \pause
  temos
  \[
    f(x, y)\evalat{y=mx}{}
    \pause
    = \frac{2x^2y}{x^4+y^2} \evalat{y=mx}{}
    \pause
    = \frac{2mx^3}{x^4+m^2x^2}
    \pause
    = \frac{2mx\,x^2}{\left(x^2+m^2\right)x^2}
    \pause
    = \frac{2mx}{x^2+m^2}
  \]

  \pause
  $f$ tende para zero quando $x$ vai para zero por todas as direções

  \pause
  Portanto,
  \[
    \lim_{(x, y)\to(0, 0)} \frac{2x^2y}{x^4+y^2} = 0
    \pause
    \qquad
    \qquad
    \text{\textcolor{red}{\textbf{Falso!}}}
  \]
\end{frame}

%------------------------------------------------------------------------------%
\begin{frame}{Exemplo \theexemplo}
  Fazendo \(y=kx^2\), para \(k\neq 0\) e \(x\neq 0\),
  \pause
  temos
  \[
    f(x, y)\evalat{y=kx^2}{}
    \pause
    = \frac{2x^2y}{x^4+y^2} \evalat{y=kx^2}{}
    \pause
    = \frac{2kx^4}{x^4+k^2x^4}
    \pause
    = \frac{2kx^4}{\left(1+k^2\right)x^4}
    \pause
    = \frac{2k}{1+k^2}
  \]

  \pause
  Para cada valor de $k$, a função $f$ se aproxima da origem com um valor diferente

  \pause
  Portanto, o limite não existe
\end{frame}

%------------------------------------------------------------------------------%
\begin{frame}{Continuidade de Compostas}
  Se $f(x, y)$ é contínua em $(x_0, y_0)$

  \pause
  e $g$ é uma função, de uma variável, contínua em $z_0 = f(x_0, y_0)$

  \pause
  então a função composta
  \(
    h = g\circ f
  \)
  \pause
  definida como
  \[
    h(x, y) = g\big(f(x, y)\big)
  \]
  \pause
  é contínua em $(x_0, y_0)$
\end{frame}

%------------------------------------------------------------------------------%
\addtocounter{exemplo}{1}
\begin{frame}{Exemplo \theexemplo}
  Mostre que
  \[
    f(x, y) = \cos\left(\frac{xy}{x^2+1}\right)
  \]
  é contínua em todos os pontos do plano
\end{frame}

%------------------------------------------------------------------------------%
\begin{frame}{Exemplo \theexemplo}
  Primeiro reconhecemos que temos a composição de funções $f=g\circ h$,
  \pause
  com
  \[
    h(x, y) = \frac{xy}{x^2+1}
    \qquad\qquad\text{e}\qquad\qquad
    g(t) = \cos(t)
  \]
  \pause
  isso é,
  \[
    f(x, y) = g\big(h(x, y)\big) = \cos\left(\frac{xy}{x^2+1}\right)
  \]
\end{frame}

%------------------------------------------------------------------------------%
\begin{frame}{Exemplo \theexemplo}
  $h$ é uma combinação algébrica de funções contínuas e nunca teremos
  uma divisão por zero, portanto, $h$ é contínua em todo o plano

  \pause
  \vspace{\baselineskip}
  $g$ é uma função contínua em toda a reta real, portanto, qualquer que seja o ponto
  $(x_0, y_0)$ a função $g$ será contínua em $h(x_0, y_0  )$

  \pause
  \vspace{\baselineskip}
  Pelo teorema da continuidade da composta, $f$ é contínua em todo o plano
\end{frame}

%------------------------------------------------------------------------------%
\begin{frame}{Teorema do Valor Extremo}
  Se \(f\colon\R^n\to\R\) é uma função \emph{contínua}

  \pause
  em uma região $R$, \emph{fechada e limitada},

  \pause
  a função assume um máximo e um mínimo absolutos m $R$
\end{frame}

%------------------------------------------------------------------------------%
\section{Lista Mínima}

%------------------------------------------------------------------------------%
\begin{frame}{Lista Mínima}
  Cálculo Vol. 2 do Thomas 12\textsuperscript{a} ed. -- Seção 14.2
  \begin{enumerate}
    \item Estudar o texto da seção
    \item Resolver os exercícios: 31-33, 36-37, 60
  \end{enumerate}

  \vfill
  Atenção: A prova é baseada no livro, não nas apresentações
\end{frame}

%------------------------------------------------------------------------------%
\end{document}
%------------------------------------------------------------------------------%
